\section*{Introduction}
\addcontentsline{toc}{section}{Introduction}
\label{sec:introduction}

Fourteen-year-old Sewell Setzer III did not encounter a foundation model in a laboratory. He encountered a deployed relationship. In \emph{Garcia v. Character Technologies, Inc.}, his mother's complaint alleged that Character.AI combined persistent personas, memory, access rules, warnings, engagement design, and safety controls in a consumer service available to minors. It alleged that Character Technologies' founders and Google made different contributions of personnel, technology, knowledge, and resources and that, after months of interaction with a character modeled on a fictional television figure, Sewell died by suicide. His mother brought wrongful-death, negligence, product-defect, failure-to-warn, deceptive-practice, and related claims. At the pleading stage, the district court treated the application as plausibly alleged to be a product and declined, on that record, to classify the generated output as protected speech. The action later settled without a finding of defect, causation, or liability.\lrfootnote{See \casecite[1165--69, 1174--84]{garcia2025}; \casecite[1]{garciaDismissal2026}; \casecite[1]{garciaLien2026}. The complaint's disputed account remains an allegation. This Article uses the published Rule 12 order for its classification and pleading analysis, not as a merits judgment.}

The case exposes a problem analytically independent of the ultimate merits result. The user sees the interface and the output. The facts ordinary law needs lie behind both: which model and application versions ran; who selected the relational objective, system instructions, memory, age gate, crisis response, and engagement metric; what evaluations, complaints, and incident signals arrived; who could update, restrict, roll back, or stop the deployment; and which enterprise retained each record. In a direct-to-consumer service, the operating company may be obvious while its deployment record remains inaccessible. In a divided system, the visible hospital, employer, agency, or application can depend on a model vendor, scoring service, infrastructure provider, and systems integrator that retain different controls. The claimant needs the configuration and control record to test a recognized claim, yet often must identify the record holder before obtaining that record. The merits standard exists; access to the facts that feed it does not.

This Article calls that sequence the \term{pre-merits attribution gap}. It comprises two related failures. The first is an \term{actor-identification failure}: an affected person cannot tell which enterprise selected or retained a control relevant to the asserted risk. The second is a \term{deployment-record failure}: even when the operator's name is visible, the versions, configuration, evaluations, warnings, incidents, and intervention authority needed to test its conduct remain internal or were never recorded. The first failure is most acute in divided deployments. The second can determine a case in any deployment.

AI-liability scholarship disagrees over the substantive rule that should govern after the relevant actor and conduct are known. Anat Lior argues that legislators should adopt strict liability for AI-inflicted damages, while Zachary Henderson rejects categorical strict liability, arguing that negligence and products liability should remain the defaults, supplemented by insurance and narrow legislation. Catherine Sharkey identifies products liability as the most promising framework for AI, while Ketan Ramakrishnan offers a qualified defense of negligence as the principal doctrinal foundation for foundation-model liability.\lrfootnote{See \artcite[91, 94--97]{liorStrictLiability2021}; \artcite[483, 515--17]{hendersonRobotLiability2026}; \artcite[240, 244--47]{sharkeyProductsLiability2024}; \artcite[960--61]{ramakrishnanTortFrontier2026}. A systematic review of 559 federal opinions likewise reports that courts predominantly resolve AI disputes through preexisting legal doctrines and that litigated harms are shaped by existing causes of action. See \bluecite[1, 8--9]{yuVisibleCourt2026}.} Those positions address the merits-side question: which rule should govern an identified actor and deployment. This Article addresses a logically prior access question: how an affected person identifies the deployment, actor, and bounded record to which any such rule would apply. The merits firewall leaves that substantive disagreement intact. Negligence cannot evaluate conduct until the relevant conduct and actor are visible. Products liability cannot identify which company selected the model harness or retained the update right. Strict liability still requires identification of the covered activity or product and its connection to the injury. A rule placing a category of loss outside liability still requires enough information to determine whether the case falls within that category.

The question presented is therefore narrow: when an adverse event affects an interest protected by specified law through a designated high-impact AI deployment, but the affected person cannot identify the operative system, control holder, or deployment record, what facts should open a bounded identity-and-event stage---and which enterprise must answer first?

The Article's contribution is institutional sequencing. Section~\ref{sec:adjacent-literature} maps neighboring proposals that begin with a substantive duty, disclosure mandate, regulator, or an access problem after an actor and artifact are identifiable; existing procedure offers discovery after a viable action reaches an identifiable target. The front door brings together four functions: a sectoral designation that identifies the legal interest; an ex ante minimum record that keeps the deployment observable; a visible-operator gateway that supplies a reachable first recipient; and a material-control test that allocates access without allocating liability. The resulting rule is narrow enough to administer and complete enough to open the merits process.

Recent work has made the adjacent evidence problem explicit on both sides of litigation's filing line. Stewart proposes a pre-litigation right to test how modified inputs would change a known algorithmic decision. Cen, Ismael, and Zheng identify the ``privatization of proof'' and structure judicial access disputes once a requester can identify an artifact and a producing party that controls it.\lrfootnote{See \bluecite[1067, 1076--79]{stewartBeyondExplanation2026}; \bluecite[7794--95, 7801--05]{cenPrivatizationProof2026}.} This Article supplies a different institutional layer: the front end that identifies the deployment and its control holders, creates a minimum deployment record before an adverse event becomes a lawsuit, and gives the visible operator a bounded first-answer duty. Counterfactual testing, ordinary discovery, and audit can then operate on a known deployment, version, and record.

The primary adoption audience is a state legislature regulating a designated sector. Congress can use the same architecture for a national market, a federal cause of action, or an interstate deployment requiring uniformity. Courts administer the closed first-answer proceeding, protect the record, and decide whether Stage Two is justified. They do not create the ex ante record duty, nonwaiver rule, tolling period, or enforcement remedy from ordinary discovery doctrine. In short, this Article proposes model sectoral legislation that courts can administer through a closed pre-merits proceeding.

This Article proposes a \term{deployment front-door rule}. A legislature first designates a deployment category and specifies the protected interests, visible operator, minimum record, retention period, trigger, and remedy. An affected person invokes Stage One by identifying an adverse event affecting an interest protected by the specified law, the covered deployment through which it occurred, and facts reasonably connecting the requested identity and event fields to that occurrence. Stage One does not require facts held exclusively in the deployment record. The enterprise that made the deployment available or incorporated it into the affected process becomes the disclosure gateway. It must identify the deployed system, provide a bounded event receipt and custodian map, preserve the corresponding record, and activate the contractual retrieval path. An identified enterprise remains for Stage Two only when it held practical authority over a risk-relevant function, the Stage One record supports that the function formed part of the event path or governed a safeguard, and the requested information bears non-de-minimally on an element, defense, or remedy.

The rule changes access to facts, not the elements that decide the case. Its merits firewall creates no merits inference. Jurisdiction, immunity, privilege, proportionality, trade-secret protection, and substantive defenses remain operative. An enterprise that completes the first-answer obligation and establishes that it lacked material control over the risk terminates its special front-door role; absent an independently sufficient claim, it returns to ordinary nonparty status. Ultimate allocation remains governed by the cause of action, contribution and indemnity law, contract, and any independently available individual claim.

The front door has an ex ante component because litigation cannot recover a record that a deployment never created. For each designated category, the statute should impose a \term{deployment-record duty}: identify the system and version; assign risk and approval authority; document material evaluations, unresolved risk decisions, changes, complaints, incidents, and interventions; identify who can restrict, roll back, or stop operation; and retain the record for a specified period. A responsibility charter is the internal instrument that satisfies that record duty: infrastructure for the front door, not a second governance product.

The scope of the required record follows practical control. This Article uses a compressed \term{deployment-control cascade} to map objectives, resources, configuration, feedback, and stopping power across the model provider, application deployer, and operating institution. The cascade is a record-specification method. It identifies evidence relevant to authorization, notice, foreseeability, feasible precaution, causal sequence, preservation, and defense; governing law supplies every legal standard. Token-level unpredictability may complicate causal proof, but it does not make organizational choices about purpose, users, data, instructions, tools, monitoring, updates, and shutdown disappear.

Three design constraints define the rule's specifications. First, ordinary pleading and discovery often work once a claimant has a proper defendant and an identifiable record holder. The gap lies in establishing those predicates and in creating the record before litigation. Second, identification and production impose real costs even when the enterprise ultimately prevails. Coverage must therefore be sectorally designated; Stage One must be closed and event specific; Stage Two must be claim specific and proportionate; and privacy, security, privilege, and trade secrets must be protected. Third, some losses connected to AI have no legal remedy. The designation therefore names the protected interest and enforcement route, pure legal defenses can be resolved before Stage Two, and the statute supplies no general right to inspect a deployment simply because an affected person invokes AI.

The Article develops one contribution through five Parts. Part I identifies the pre-merits attribution gap, distinguishes adjacent evidence-access scholarship, and compresses the deployment-control cascade into a tool for specifying the missing record. Part II states the front-door rule, separates Stage One from Stage Two, defines material control, supplies model statutory text, and coordinates the substantive access right with federal procedure. Part III applies the rule to two complementary settings: \emph{Garcia}, where the visible operator makes production the central problem, and \emph{Mobley}, where feature identity, data hosting, legal control, and validation divide across employers and a vendor. Part IV develops a minimized ex ante record infrastructure and compares existing safety, information-demand, and private-access regimes. Part V defines four boundaries: existing legal interests; constitutional and confidentiality protections; arbitration, cross-border enforcement, and preemption; and classification neutrality. The Conclusion returns to a single proposition: substantive law should decide AI cases, not an information asymmetry that prevents substantive law from starting.
